\section{Linear Regression}

\subsection{Univariate Linear Regression}

\subsubsection{Model}

Univariate linear regression is a linear relationship 
between an independent variable $x \in \real$, 
also called the 
{\em explanatory variable}, and a dependent variable
$y \in \real$, also called a {\em response variable}.

The {\bf model function} is
\be
y = w x + b
\label{eq:lin_reg_model}
\ee
where
\begin{itemize}
\item $w \in \real$ is the slope (also called {\bf weight}) 
acting on the independent variable, and
\item $b \in \real$ is the intercept (also called {\bf bias}).
\end{itemize}
The two real-valued numbers ($w, b$) are called
{\bf model parameters}.

\subsubsection{Data}

Assume $m$ pair of observations ($x_j, y_j$), $j=1\ldots m$, 
have been made.  
Assume equal likelihood of each observation.
Given the observations, the objective
of univariate linear regression is to 
\begin{enumerate}
\item find model parameters $(w, b)$ that {\bf best fit} the 
observed data, and 
\item quantify the {\bf correlation coefficient}, 
which is an objective measure of 
how well the best fit model describes the underlying 
observed data.  
\end{enumerate}

\subsubsection{Error Function}
 
For each observation ($x_j, y_j$) pair, there is an error
$e_j(w,b)$ that is defined as the difference between the 
actual observed response $y_j$ and the model's 
predicted response $w x_j + b$, defined as
\be
 e_j(w,b) \defe y_j - \left(w x_j + b \right).
\label{eq:L1_norm_lin_reg}
\ee
The error accumulated for all $m$ observations can be 
written as
\be
\parallel \ve{e} (w, b)
\parallel_1 \;  \defe \sum_{j=1}^m \abs{e_j(w,b)}.
\ee
This is the L1-norm of the error vector $\ve{e}
= [e_1, e_2, \ldots, e_m] \in \real^m$.
See Section~\ref{sec:pnorm} for more details on p-norms.

The model parameters $(w, b)$ are said to be best fit
when the error is minimized.  While the L1-norm error
in~(\ref{eq:L1_norm_lin_reg}) may be used as the objective
error function, the L2-norm is frequently used as an 
alternative to the L1-norm.

For this model, the L2-error norm is defined as
\be
\parallel \ve{e} (w, b)
\parallel_2 \;  \defe \sum_{j=1}^m \left(e_j(w,b)\right)^2 
= \sum_{j=1}^m \left(
 y_i - w x_i - b
\right)^2.
\label{eq:L2_norm_lin_reg}
\ee
Using the notation $E_2 \in \real$ to denote the L2-error norm
in~(\ref{eq:L2_norm_lin_reg}),
we expand and simplify the squared term,
\be
E_2(w,b) = 
\sum_{j=1}^m y_j^2 - 2 y_j w x_j - 2 y_j b + w^2 x_j^2 
+ 2 w x_j b + b^2.
\label{eq:E_2_w_b}
\ee
Since $E_2$ is the summation of squared quantities, 
$E_2(w,b) \geq 0\;\forall\; w \in \real, b \in \real$.
$E_2(w,b)$ 
behaves as a quadratic, scalar function.
Values of $(w, b)$ that minimize the error $E_2$ are 
found by taking partial derivatives of the error function 
with respect to the parameters $w$ and $b$:
\be
\frac{\partial E_2}{\partial w} = 
\sum_{j=1}^m -2 y_j x_j + 2 w x_j^2 + 2 x_j b
\label{eq:partial_E_w}
\ee
and
\be
\frac{\partial E_2}{\partial b} = 
\sum_{j=1}^m -2 y_j + 2 w x_j + 2 b.
\label{eq:partial_E_b}
\ee
Setting~(\ref{eq:partial_E_b}) equal to zero, and 
using the definitions of $\mu_x$ and $\mu_y$
from~(\ref{eq:simple_average}),
\be
b = \mu_y - w \mu_x \; \mbox{ to minimize } E_2.
\label{eq:b_for_E_2_min}
\ee
Substituting
(\ref{eq:b_for_E_2_min}) into 
(\ref{eq:E_2_w_b}) for values of $b$ that minimize
$E_2$ allows
(\ref{eq:E_2_w_b}) to be written as a function of 
$w$ alone, 
\begin{align}
E_2(w) = \sum_{j=1}^m & 
\; y_j^2 - 2 y_j w x_j - 2 y_j (\mu_y - w \mu_x) + 
w^2 x_j^2 + \nonumber \\
& 2 w x_j (\mu_y - w \mu_x) + 
(\mu_y - w \mu_x)^2.
\end{align}
Then, 
\begin{align}
\frac{d E_2}{dw} = 
\sum_{j=1}^m & 
\; -2 y_j x_j + 2 y_j \mu_x + 2 w x_j^2 + \nonumber \\
& 2 x_j \mu_y - 4 w x_j \mu_x +
2(\mu_y - w \mu_z)(-\mu_x)  
\overset{\mbox{\tiny set}}{=} 0, 
\end{align}
or
\begin{align}
0 & = \sum_{j=1}^m 
-y_j x_j + y_j \mu_x + w x_j^2 + 
 x_j \mu_y - 2 w x_j \mu_x + w \mu_x \mu_x - \mu_x \mu_y, \\
 & = \sum_{j=1}^m
 -\left(x_j - \mu_x\right)
  \left(y_j - \mu_y\right) 
  + 
  w \left(x_j - \mu_x\right)^2.
\end{align}
Solving for $w$
\be
w = \frac{\dst\sum_{j=1}^m 
  \left(x_j - \mu_x\right)
  \left(y_j - \mu_y\right)}
  {\dst\sum_{j=1}^m
  \left(x_j - \mu_x\right)^2}.
\ee
Multiplying both the numerator and the denominator
by $\frac{1}{m}$ gives
\be
w = \frac{\cov(X,Y)}{\var(X)}.
\ee
Thus the first objective of univariate linear regression, 
find model parameters ($w, b$) the best fit to the 
observed data (by minimizing
the error between the model and the data) has been accomplished.
The second and final objective, to quantify the quality of
the best fit, follows.

Using the sample correlation coefficient
$\rho_{X,Y}$, the sample (uncorrected) standard 
deviation $\sigma_x$ and $\sigma_y$, the sample
variance and covariance,
\begin{align}
\var(X) & = \sigma_x^2, \\
\cov(X,Y) &= \rho_{X,Y} \sigma_x \sigma_y,
\end{align}
then
\be
w = \frac{\rho_{X,Y} \sigma_y}{\sigma_x} \implies
\rho_{X,Y} = \frac{w \; \sigma_x}{\sigma_y}.
\ee

\begin{Hresult}
As long as the univariate linear regression model includes the 
intercept (bias) term $b$ (\ie the 
model is not forced to go through the origin by virtue of 
$b \overset{\mbox{\tiny set}}{=} 0$.), then 
\begin{enumerate}
\item the sum of the errors is zero
\be
\sum_{j=1}^m e_j = 0, \;\; \mbox{and}
\ee
\item the model regression line goes through the center of 
mass point ($\mu_x, \mu_y$).
\end{enumerate}
\end{Hresult}

\begin{Hproof}
For part 1, 
start from the definition~(\ref{eq:L1_norm_lin_reg}),
\be
e_j = y_j - w x_j - b,
\ee
add up all the errors, $j=1\ldots m$, and divide by $m$,
\be
\frac{1}{m} \sum_{j=1}^m e_j  = 
  \frac{1}{m} \sum_{j=1}^m y_j - w x_j - b.
  \ee
Substitute the definitions of $\mu_x$ and $\mu_y$
from~(\ref{eq:simple_average})
  \be
  = \mu_y - w \mu_x - b \ee
and finally, substitute the model 
equation~(\ref{eq:lin_reg_model}) with values of
$y \mapsto \mu_y$ and $x \mapsto \mu_x$,
\be = \mu_y - w \mu_x - \mu_y + w \mu_x = 0 \ee
which is the result for part 1.

For part 2, show that the error at the point ($\mu_x, \mu_y$) 
is zero, 
\be
\left.e_j\right|_{(\bar{x},\bar{y})}
 = \left.y_j\right|_{\bar{y}} - 
(w \left.x_j\right|_{\bar{x}} + b) = \frac{1}{m}\sum_{j=1}^m y_j 
 - \left( 
 \frac{w}{m}\sum_{j=1}^m x_j + b
 \right).
\ee
Rearranging, 
\be
\left.e_j\right|_{(\mu_x,\mu_y)}
 = 
 \frac{1}{m}\sum_{j=1}^m 
 \left(
 y_j - (w x_j + b)
 \right)
 = 
 \frac{1}{m}\sum_{j=1}^m 
 e_j = 0
\ee
by the result in part 1.  This 
shows that the point ($\mu_x, \mu_y$) is on the 
line described by the model in~(\ref{eq:lin_reg_model}).
\end{Hproof}

